% 可替换原5.3节；引用键 p1-toussaint1983 沿用主论文。
\subsection{定位区域直径的旋转卡壳求解}
\label{sec:p1-calipers}

\subsubsection{将区域直径转化为顶点间的最大距离}

上一节已获得定位区域的有序顶点 $V_0,V_1,\ldots,V_{h-1}$。以下将非空有界定位区域简记为 $P$。题目要求计算区域内任意两点之间距离的最大值，即
\begin{equation}
    D(P)=\max_{X,Y\in P}\|X-Y\|_2.
\end{equation}
这里的 $X$ 和 $Y$ 可以位于区域内部、边界或顶点，并不预先限定为多边形顶点。直接搜索区域内的全部点不可行，因此需要利用凸性将搜索范围缩小到有限个顶点。

由于 $P$ 是凸多边形，任意 $X,Y\in P$ 都可以表示为各顶点的加权平均：
\begin{equation}
    X=\sum_{i=0}^{h-1}\alpha_iV_i,\qquad
    Y=\sum_{j=0}^{h-1}\beta_jV_j,
\end{equation}
其中 $\alpha_i,\beta_j\ge0$，且 $\sum_i\alpha_i=\sum_j\beta_j=1$。因此，
\begin{equation}
\begin{aligned}
    \|X-Y\|_2
    &=\left\|\sum_{i,j}\alpha_i\beta_j(V_i-V_j)\right\|_2\\
    &\le\sum_{i,j}\alpha_i\beta_j\|V_i-V_j\|_2\\
    &\le\max_{i,j}\|V_i-V_j\|_2.
\end{aligned}
\end{equation}
也就是说，区域内任意两点的距离都不会超过最远顶点对的距离。另一方面，顶点本身也属于区域，故区域直径恰好为
\begin{equation}\label{eq:p1-vertexdiam}
    D(P)=\max_{0\le i,j<h}\|V_i-V_j\|_2.
\end{equation}
于是，只需在顶点之间寻找最远点对，便可得到整个连续区域的直径。穷举全部 $\binom h2$ 个顶点对需要 $O(h^2)$ 时间，下面利用旋转卡壳进一步减少搜索量。

\subsubsection{平行支撑线与候选点对}

设一条直线与凸多边形接触，且整个多边形都位于它的同一侧，则称这条直线为支撑线。用两条相对的平行支撑线夹住多边形时，两条线上分别接触的顶点构成对踵点对。旋转卡壳的基本思想是让这两条支撑线围绕多边形转动，并随转动更新接触顶点\cite{p1-toussaint1983}。

之所以可以重点搜索对踵点对，是因为最远顶点对必然具有这样的支撑关系。设最远顶点为 $A,B$，记 $D=\|B-A\|_2>0$，取单位方向向量 $u=(B-A)/D$。对任意 $X\in P$，都有
\begin{equation}\label{eq:p1-support}
    0\le(X-A)\cdot u\le D.
\end{equation}
如果左侧不等式不成立，则 $X$ 沿 $AB$ 方向位于 $A$ 的外侧，其到 $B$ 的距离将超过 $D$；如果右侧不等式不成立，则 $X$ 位于 $B$ 的外侧，其到 $A$ 的距离将超过 $D$。这两种情况都与 $A,B$ 为最远点对相矛盾。因此，整个多边形必位于过 $A,B$ 且垂直于 $AB$ 的两条直线之间，这两条直线正是相对的平行支撑线。

这一性质说明，只要完整处理旋转过程中可能实现最远距离的支撑点组合，就不会遗漏区域直径。

\subsubsection{利用三角形面积推进支撑点}

对逆时针排列的顶点，固定当前边 $V_iV_{i+1}$，定义
\begin{equation}\label{eq:p1-area}
    F_i(j)=(V_{i+1}-V_i)\times(V_j-V_i),
\end{equation}
其中下标按模 $h$ 循环计数。$F_i(j)$ 是三角形 $V_iV_{i+1}V_j$ 面积的两倍，且满足
\begin{equation}
    F_i(j)=\|V_{i+1}-V_i\|_2\,operatorname{dist}(V_j,L_i),
\end{equation}
其中 $L_i$ 是当前边所在的直线。由于边长固定，面积最大的顶点也是距离 $L_i$ 最远的顶点，即对侧支撑线的接触点。

在计算过程中，比较 $F_i(j+1)$ 与 $F_i(j)$：如果前者较大，就把对侧顶点从 $V_j$ 推进到 $V_{j+1}$，直到面积不再增加。这里最大化的是点到当前边所在直线的距离，并不直接等于区域直径；仍需计算相应顶点对的欧氏距离。

凸多边形的边方向沿边界有序变化，对应的支撑位置也按边界顺序变化。因此，当当前边移到下一条边时，对侧支撑点可以从已有位置继续向前搜索，无须回到起点。正是这种单向推进性质，避免了对每条边重新遍历全部顶点。

若 $F_i(j+1)=F_i(j)$ 且两者均达到最大值，则对侧支撑线同时经过两个相邻顶点，即与一条边重合。此时需要把两个对侧端点均纳入候选，完整处理平行边情形。

\subsubsection{算法流程与结果的可靠性}

综合上述关系，采用以下扫描步骤。下述比较首先按精确运算描述，数值实现中的容差另行讨论。
\begin{enumerate}[label=\arabic*)]
    \item 若区域为单点，返回 $D=0$；若区域为线段，返回两端点距离。对 $h\ge3$ 的凸多边形，初始化对侧指针 $j=1$ 和最大距离平方 $M=0$。
    \item 依次处理各边 $V_iV_{i+1}$。当 $F_i(j+1)>F_i(j)$ 时推进 $j$，直至到达当前边的对侧支撑位置。
    \item 分别计算 $\|V_i-V_j\|_2^2$ 和 $\|V_{i+1}-V_j\|_2^2$，与 $M$ 比较并保存较大值及其端点。
    \item 若对侧出现等面积的相邻顶点，则进一步检查 $\|V_i-V_{j+1}\|_2^2$ 和 $\|V_{i+1}-V_{j+1}\|_2^2$，更新最优结果。
    \item 扫描全部边后，输出 $D=\sqrt M$ 以及对应最远点对。
\end{enumerate}

上述扫描覆盖最远点对的原因可以结合支撑线理解：从任一最远点对对应的平行支撑线出发，在保持这两个顶点接触的条件下转动支撑线，直至其中一条线首次与多边形的一条边重合；若初始已经重合，则无需转动。此时最远点对必属于该边端点与对侧支撑端点的组合。算法依次处理每条边，并在平行边情形下检查全部端点组合，因此会检查到实现直径的点对。同时，所检查的点都属于定位区域，其距离不可能超过区域直径，故扫描所得最大值与式\eqref{eq:p1-vertexdiam}一致。

整个过程的适用性取决于交会区域非空、有界，以及顶点序列和支撑事件处理的完整性，不要求检测点等距或均匀分布，也不要求定位区域具有固定的边数。边指针遍历 $h$ 条边，对侧指针仅沿边界推进常数圈，故已有凸包上的直径扫描时间复杂度为 $O(h)$。
